% Related work for the BISON manuscript.
% Author: fangtoast
% Email: fangtoast@foxmail.com
% Copyright (c) 2026 fangtoast. All rights reserved.

\section{Related Work}

\subsection{From Solver-Based Photonic Design to Learned Search}

MPB solves the frequency-domain Maxwell eigenproblem in a plane-wave basis
~\cite{johnson2001mpb}. Topology optimization couples such physical models to
gradient-based search over many material degrees of freedom
~\cite{molesky2018inverse}. This approach has produced low-loss photonic-crystal
bends~\cite{jensen2004systematic}, nanophotonic components, and
dispersion-shaped or manufacturing-robust waveguides
~\cite{jensen2011topology,wang2011robust}. Grating couplers, wavelength
demultiplexers, and fabrication-constrained devices further demonstrate the
scope of solver-coupled optimization
~\cite{piggott2014grating,piggott2015demultiplexer,piggott2017fabrication}.
Its main computational cost is the repeated solution of a new physical problem
for each objective and design iterate.

Learned surrogates shift part of this cost to dataset generation and training.
They have accelerated nanoparticle simulation, power-splitter design, and
global design-space exploration
~\cite{peurifoy2018nanophotonic,tahersima2019splitters,melati2019mapping}.
Tandem networks address inverse non-uniqueness by training through a frozen
forward surrogate~\cite{liu2018tandem}, whereas conditional generators and
global-optimization networks search over multiple candidates
~\cite{jiang2019glonet}. Across direct, iterative, and generative approaches,
the recurring concerns are non-uniqueness, data coverage, surrogate bias, and
physical validation~\cite{so2020inverse,ma2021deep,wiecha2021deep,
jiang2021evaluation}. Budget-aware evaluation is therefore essential
~\cite{ren2020benchmarking}. Our protocol fixes both the \(K=128\) candidate
pool and the Top-5 MPB replay budget.

\subsection{Neural Operators and Fourier-Space Inductive Bias}

DeepONet represents nonlinear operators with branch and trunk networks
~\cite{lu2021deeponet}. FNO alternates pointwise transformations with global
convolutions parameterized by learned Fourier modes~\cite{li2021fno}. Both fit
within the broader framework of learning maps between function spaces
~\cite{kovachki2023neuraloperator}. Subsequent variants extend Fourier mixing
or geometry handling. AFNO adds adaptive sparsification and channel structure
~\cite{guibas2022afno}, Geo-FNO deforms irregular domains to a uniform latent
grid~\cite{li2023geofno}, and NUNO decomposes non-uniform domains into
grid-aligned subdomains~\cite{liu2023nuno}.

BISON addresses a narrower setting with a stronger structural prior: a fixed
periodic Cartesian cell. It specializes Fourier mixing to two configured mode
rectangles and jointly predicts an ordered band matrix. Unlike
physics-informed neural operators, which can incorporate equation residuals
~\cite{li2021pino}, BISON is trained on paired permittivity and eigenspectrum
data. Neural Inverse Operators instead learn operator compositions for PDE
inverse problems~\cite{molinaro2023nio}. In this paper, ``operator'' denotes the
learned functional map from a permittivity field to a finite sampled
eigenspectrum.

\subsection{Learned Electromagnetics and Photonic Band Surrogates}

Learned electromagnetic models differ in both prediction target and validation
evidence. MaxwellNet incorporates Maxwell-equation residuals into field
prediction~\cite{lim2022maxwellnet}. NeurOLight and PACE use neural operators
for parametric photonic-device field simulation
~\cite{gu2022neurolight,zhu2024pace}, and an operator surrogate has supported
free-form electromagnetic inverse design~\cite{augenstein2023emno}. Field
solutions, port responses, and ordered eigenfrequency matrices nevertheless
define different learning tasks and require different evaluation protocols.

Christensen et al. predict complete band structures and generate band-gap
candidates for two-dimensional photonic crystals
~\cite{christensen2020predictive}. POD--DeepONet similarly studies forward and
inverse band design on binary, symmetric two-dimensional cells
~\cite{wang2026poddeeponet}. Direct accuracy comparisons are not warranted
because the datasets, discretizations, targets, and replay protocols differ.
BISON instead combines joint full-band prediction with a routing-free latent
generator, an independent query audit, and equal candidate and MPB budgets.

\begin{table*}[t]
  \caption{Positioning of BISON among photonic inverse-design methods.
  ``Physical check'' summarizes the typical evidence type; solvers and budgets
  differ across the cited studies.}
  \label{tab:positioning}
  \centering
  \small
  \begin{tabular}{@{}p{0.18\textwidth}p{0.20\textwidth}p{0.20\textwidth}p{0.17\textwidth}p{0.17\textwidth}@{}}
    \toprule
    Family & Learned or optimized object & Main advantage & Main risk or cost & Physical check \\
    \midrule
    Topology/adjoint design & One structure per objective & Direct solver gradients & Many iterative solves & In-loop or final solver \\
    Direct/tandem learning & Target response $\rightarrow$ structure & Fast amortized inference & One-to-many ambiguity; surrogate bias & Often post-hoc simulation \\
    Generative/search learning & Candidate distribution or latent search & Multiple feasible designs & Budget-sensitive ranking & Candidate-dependent \\
    Photonic neural operators & Geometry/conditions $\rightarrow$ fields or response & Global functional mapping & Distribution shift & Field/response solver \\
    BISON workflow & Periodic cell $\rightarrow$ full bands; latent candidates & Joint band ranking with explicit budgets & Fixed task family; surrogate error & Top-5 MPB replay per target \\
    \bottomrule
  \end{tabular}
\end{table*}
